In this section, we describe the empirical evaluation datasets, baselines, ablations, and metrics, and then discuss the main numerical results.

\subsection{Datasets}
\label{sec:data}
Evaluation data were obtained from the open-access PK-DB repository, which aggregates pharmacokinetic measurements from clinical and pre-clinical studies~\citep{grzegorzewski2021pkdb}.
%
From PK-DB, we curated plasma concentration--time measurements from a Phase~I study in healthy volunteers investigating the safety and dosing of a combination of drugs and their metabolites~\citep{lenuzza2016safety}.
%
In total, this dataset comprises 18 compounds, including parent drugs and primary metabolites, with 2019 valid concentration--time observations.
%
Sampling frequencies vary substantially both within and across compounds, resulting in sparse and irregularly sampled time series.
%
Each compound is measured in 10 individuals, with up to 15 observations per individual.
%
We also included two additional pharmacokinetic datasets on indometacin and theophylline (\texttt{Indometh} and \texttt{Theoph} from the R \texttt{datasets} package)~\citep{Rsoftware}.
%
Together, these 20 empirical compound-level datasets reflect the small-cohort and sparsely sampled regime targeted by PFF.

\subsection{Baselines and Ablations}
\label{sec:baselines}

We compare Prior-Fitted Functional Flows (PFF) against one classical pharmacometric baseline and three neural baselines: nonlinear mixed-effects modeling (NLME), NODE-PK, Transformer, and \texttt{AICMET}.
%
The NLME baseline is fitted study-by-study for each parent compound using standard estimation procedures. We fit one- and two-compartment models with first-order absorption, inter-individual variability, and combined residual error models, and select the final model by AIC. Metabolites are excluded from the NLME comparison because they lack explicit dosing inputs.
%
In contrast, all neural models are pretrained on synthetic PK studies and evaluated in context on empirical studies without dataset-specific fine-tuning.
%
NODE-PK extends latent neural ODEs~\citep{rubanova2019latent} with explicit dosing information: an RNN encoder produces an individual latent initial condition, which is then evolved by neural ODE dynamics~\citep{lu2021deep}. NODE-PK treats individuals independently and does not aggregate study-level information.
%
The Transformer baseline replaces NODE-PK's neural ODE dynamics with a functional-attention decoder, but does not use the full hierarchical mixed-effects representation of \texttt{AICMET}.
%
\texttt{AICMET}~\citep{marin2025amortized} is a hierarchical neural-process baseline that encodes study-level and individual-level effects through latent mixed/random-effect representations and decodes concentration functions with functional attention.
%
Relative to \texttt{AICMET}, PFF changes both the inferential object and the vector-field architecture: it learns a direct transport between measures over concentration functions, and parameterizes the transport vector field with continuum-attention neural operators rather than an AICMET-style functional-attention decoder.

For PFF, we additionally evaluate two ablations that isolate central design choices.
%
First, \texttt{PFF-GP-prior} replaces the GP posterior reference measure used for partially observed individuals by an unconditional GP prior, thereby isolating the effect of prefix-conditioned initialization.
%
Second, \texttt{PFF-std-attn} replaces continuum attention in the vector-field architecture by standard attention, recovering an AICMET-style functional-attention mechanism while keeping the flow-matching objective fixed.


\subsection{Metrics}

\paragraph{Forecasting evaluation.}
Forecasting performance is assessed in a leave-one-individual-out setting.
%
For subject $i$, predictions are conditioned on the study context $\mathcal{S}$, dosing inputs $\mathbf{u}^i$, and the first four PK samples $\mathcal{D}^i_{\mathrm{C}}$.
%
The resulting posterior predictive trajectories are evaluated on the remaining samples $\mathcal{D}^i_{\mathrm{T}}$ using log-RMSE.

\paragraph{Generative evaluation.}
We evaluate time-series sample quality in two complementary settings: an \emph{empirical} setting and a \emph{synthetic} setting.
%
In both cases, the goal is to compare a reference set of trajectories against a corresponding set of trajectories generated by the model under the same study-level conditioning, but the construction of the reference set differs across settings.
%
For the empirical evaluation, we apply a leave-one-out protocol across all empirical studies and drugs.
%
In each study, one individual is removed from the study context, the model is conditioned on the remaining individuals, and one trajectory is sampled for the held-out individual.
%
Repeating this procedure over all individuals in all studies yields a generated empirical set whose cardinality matches that of the empirical dataset, since each empirical individual is held out and resampled exactly once.
%
The generated trajectories are then pooled across studies and compared against the pooled empirical trajectories from the original dataset.
%
For the synthetic evaluation, no leave-one-out construction is needed.
%
Instead, each synthetic study provides a study context together with a reference population of target individuals sampled directly from the simulator.
%
Conditioned on the same context, the model generates a matching synthetic population, and the generated trajectories are compared directly against the corresponding oracle trajectories.
%
After constructing these empirical and synthetic comparison pairs, we quantify discrepancy using two complementary two-sample metrics: (i) the area under the ROC curve (AUC) of a classifier trained to distinguish reference from generated trajectories, where values closer to $0.5$ indicate better sample fidelity; and (ii) a maximum mean discrepancy (MMD) computed with the signature kernel of \citet{kiraly2019kernels}, which measures distributional discrepancy between sets of irregular trajectories in path space. Since we use an unbiased finite-sample estimator, small negative MMD estimates can occur; values closer to zero indicate better agreement.
%
We additionally assess generative behavior using visual predictive checks (VPCs), comparing simulation percentiles against observed concentrations~\citep{holford2005}. Full implementation details are provided in the Supplementary Material.

\begin{table}[t]
\centering
\caption{\textbf{Sample fidelity scores.} AUC values closer to $0.5$ and MMD values closer to zero indicate better agreement between generated and reference trajectories. Small negative MMD estimates can occur due to the unbiased finite-sample estimator.}
\label{tab:sample_fidelity}
\scriptsize
\begin{tabular}{lcc}
\toprule
\textbf{Model} & \textbf{AUC} & \textbf{MMD} \\
\midrule
\multicolumn{3}{c}{\textbf{Empirical}} \\
\cmidrule(lr){1-3}
\texttt{AICMET} & 0.577 & -0.001 \\
\texttt{PFF} & \textbf{0.549} & \textbf{0.0009} \\
\addlinespace
\multicolumn{3}{c}{\textbf{Synthetic}} \\
\cmidrule(lr){1-3}
\texttt{AICMET} & 0.541 & 0.002 \\
\texttt{PFF} & \textbf{0.508} & \textbf{0.001} \\
\bottomrule
\end{tabular}
\end{table}

\begin{table}[t]
\centering
\caption{\textbf{Ablation study on empirical sample fidelity.} We compare PFF to two ablations under the same empirical leave-one-individual-out protocol. AUC values closer to $0.5$ and MMD values closer to zero indicate better agreement with empirical trajectories.}
\label{tab:pff_ablation}
\scriptsize
\begin{tabular}{lcc}
\toprule
\textbf{Model} & \textbf{AUC} & \textbf{MMD} \\
\midrule
\texttt{PFF} & \textbf{0.549} & \textbf{0.0009} \\
\texttt{PFF-GP-prior} & 0.555 & 0.0037 \\
\texttt{PFF-std-attn} & 0.576 & 0.0014 \\
\texttt{AICMET} & 0.577 & -0.001 \\
\bottomrule
\end{tabular}
\end{table}

\subsection{Discussion of Numerical Results}
\label{sec:numerical_results}

Table~\ref{tab:main_results} summarizes the quantitative forecasting performance.
%
\texttt{PFF} achieves the lowest or competitive log-RMSE across the evaluated compounds in the majority of cases, indicating that the learned posterior predictive operator transfers effectively from simulated pretraining studies to sparse empirical PK datasets.
%
This predictive behavior is illustrated in Figure~\ref{fig:prediction_row}, which shows inferred future trajectories for partially observed individuals.
%
The model does not merely interpolate observed prefixes: it uses the cohort-level study context together with the individual prefix to infer a distribution over plausible future concentration functions.

Table~\ref{tab:sample_fidelity} evaluates zero-shot population synthesis.
%
Compared with \texttt{AICMET}, \texttt{PFF} yields an AUC closer to the ideal value of $0.5$ in both the empirical and synthetic settings, indicating that trajectories generated by PFF are harder to distinguish from reference trajectories.
%
The synthetic comparison shows the same pattern under simulator-controlled ground truth, while the empirical leave-one-individual-out comparison demonstrates that the improvement carries over to real PK studies.
%
Figure~\ref{fig:vpc_only} provides the corresponding visual predictive checks: the simulated percentiles align closely with the empirical concentration distributions across study cohorts, supporting the use of PFF for zero-shot population-level generation.

Table~\ref{tab:pff_ablation} isolates two main design choices.
%
Replacing the GP posterior reference used for forecasting with a flat GP prior degrades sample fidelity, especially in MMD, indicating that prefix-conditioned reference measures provide a useful inductive bias before transport.
%
Replacing continuum attention by standard attention has a larger effect on AUC and makes \texttt{PFF-std-attn} nearly indistinguishable from \texttt{AICMET}.
%
This suggests that the continuum/operator-attention parameterization is a key empirical driver of the improved generative fidelity on irregular PK trajectories.
%
Thus, the gain should not be attributed to the function-space flow-matching formulation alone, but to the combination of direct posterior-predictive transport in function space, GP posterior references for partially observed individuals, and continuum-attention neural operators for irregular observation grids.

Finally, PFF shifts computational cost from repeated study-specific optimization to one up-front pretraining stage.
%
At deployment time, the pretrained model can be applied off-the-shelf to new studies without retraining, whereas classical NLME workflows often require iterative model development, model comparison, and manual diagnostic inspection.
%
This makes PFF particularly relevant in early-development settings where datasets are small, irregular, and repeatedly encountered across compounds.
%
A detailed account of hyperparameters, inference times, calibration diagnostics, statistical tests, and failure cases is provided in the Supplementary Material.
